\clearpage\fancyhead[L]{\small KOMET.}\hypertarget{native-387-lower}{}\label{native-387-lower}
\begin{center}\textbf{KOMET.}\par\vspace{4pt}\hrule\vspace{4pt}\fontsize{10}{12}\selectfont Göttingische gelehrte Anzeigen. Stück 55. Seite 537\;..\;538. 8. April 1815.\par\vspace{4pt}\hrule\vspace{4pt}\end{center}

\par\indent Vom 13. bis zum 31. März gab es hier nur vier klare Abende, an denen der neue Komet beobachtet werden konnte. Hier folgen daher alle Beobachtungen, die an der hiesigen Sternwarte bisher gewonnen wurden:

\begin{center}\fontsize{9}{11}\selectfont\setlength{\tabcolsep}{3pt}\renewcommand{\arraystretch}{1.15}\begin{tabular}{l|ccc|ccc|ccc}
1815 & \multicolumn{3}{c|}{Mittlere Zeit} & \multicolumn{3}{c|}{\shortstack{Scheinbare\\Rektaszension}} & \multicolumn{3}{c}{\shortstack{Scheinbare\\Deklination}} \\

\hline
März 20 & $10^{\mathrm u}$ & $33^{\mathrm m}$ & $6^{\mathrm s}$ & $54^\circ$ & $7'$ & $1''$ & $39^\circ$ & $7'$ & $47''$ N. \\
21 & 10 & 11 & 37 & 54 & 33 & 21 & 39 & 36 & 57 \\
21 & 10 & 57 & 0 & 54 & 34 & 21 & 39 & 38 & 8 \\
25 & 9 & 46 & 38 & 56 & 28 & 50 & 41 & 38 & 5 \\
30 & 9 & 50 & 57 & 59 & 13 & 3 & 44 & 10 & 27 \\
\end{tabular}\end{center}

\par\indent Die zweite Beobachtung vom 21. März stammt von Herrn \textsc{Enke}, dessen Geschick in der praktischen Astronomie diese Blätter schon öfter belegt haben. Weitere Beobachtungen von anderen Orten sind noch nicht eingegangen.

\Needspace{19\baselineskip}\par\indent Professor \textsc{Gauss} hat aus diesen allerdings spärlichen Beobachtungen die Bahn dieses Kometen vorläufig näherungsweise bestimmt:

\begin{center}\fontsize{10}{12}\selectfont\setlength{\tabcolsep}{3pt}\renewcommand{\arraystretch}{1.15}\begin{tabular}{l}
Durchgang durch das Perihel\\
\quad Mittlere Zeit in Göttingen, 24. April 1815. $16^{\mathrm u}\;37^{\mathrm m}\;34^{\mathrm s}$\\
\end{tabular}\end{center}
\begin{center}\fontsize{10}{12}\selectfont\setlength{\tabcolsep}{3pt}\renewcommand{\arraystretch}{1.15}\begin{tabular}{lcr}
Länge des Perihels & $.\;.\;.\;.\;.$ & $246^\circ\;7'\;2''$\\
Aufsteigender Knoten & $.\;.\;.\;.\;.\;.$ & $82\;43\;4$\\
Neigung der Bahn & $.\;.\;.\;.\;.\;.\;.$ & $45^\circ\;8'\;55''$\\
Kleinster Abstand von der Sonne & $.\;.$ & $1.24738$\\
\end{tabular}\end{center}
\noindent Bewegung rechtläufig.

\SourceLeaf{388}\par\indent Aus dieser Bestimmung lässt sich schon voraussehen, dass der Komet bis gegen Ende April noch etwas heller werden wird. Diese Zunahme ist aber sehr gering. Er wird für das bloße Auge unsichtbar bleiben und daher \emph{nur für die Astronomen} interessant sein. Er wird durch den Fahrmann und den Luchs in den Großen Bären ziehen und sich dort, weil sein Licht immer weiter abnimmt, im Juni dem Auge entziehen.
